\thispagestyle{plain}

\section*{Abstract}
Fast radio bursts (FRBs) are millisecond radio transients of extragalactic origin whose
progenitor mechanisms remain unknown. Since FRBs at cosmological distances trace the large-scale
structure, their angular clustering encodes the host-population bias, which depends
on the progenitor channel. This thesis investigates
how well this bias can be constrained by cross-correlating FRBs with the galaxy sample of
KiDS-Legacy, the final data release of the Kilo-Degree Survey. Two progenitor scenarios, young
magnetars and neutron star mergers formed through delayed channels, are each described by a
redshift-dependent linear bias and evaluated for a shallow, CHIME-like FRB survey and a deep FRB
survey representative of the Square Kilometre Array. The FRB auto-correlation alone is
shot-noise dominated and does not constrain the bias. Adding the galaxy sample reduces
the marginalized uncertainties by factors of $17$ to $51$ and improves the figure of merit by factors
of $61$ to $449$, despite a sky area roughly $28$ times smaller. With the current KiDS footprint,
however, the two scenarios cannot yet be distinguished. A forecast for the Legacy Survey of
Space and Time (LSST) shows that upcoming galaxy surveys will substantially improve the prospects of
discriminating between a prompt and a delayed origin of FRBs.

\section*{Kurzfassung}
\begin{foreignlanguage}{ngerman}
Schnelle Radioblitze (engl. Fast Radio Bursts, FRBs) sind extragalaktische Radiotransienten
mit Millisekundendauer, deren Ursprung bislang unbekannt ist. Da FRBs mit
kosmologischen Entfernungen der großskaligen Struktur folgen, enthält ihre Winkelkorrelation
Informationen über den Bias der Wirtspopulation, der vom Entstehungskanal abhängt. In dieser
Arbeit wird eine Vorhersage-Pipeline entwickelt, um abzuschätzen, wie gut sich dieser Bias durch
Kreuzkorrelation von FRBs mit der Galaxienstichprobe von
KiDS-Legacy, dem finalen Datensatz des Kilo-Degree Survey, einschränken lässt. Zwei
Entstehungsszenarien, junge Magnetare und Verschmelzungen von Neutronensternen aus verzögerten Kanälen, werden jeweils
durch einen rotverschiebungsabhängigen linearen Bias beschrieben und für eine flache,
CHIME-ähnliche sowie eine tiefe FRB-Durchmusterung nach dem Vorbild des Square Kilometre Array untersucht.
 Die FRB-Autokorrelation allein
ist vom Schrotrauschen dominiert und schränkt den Bias nicht ein. Durch die Hinzunahme der
Galaxien sinken die marginalisierten Unsicherheiten um Faktoren von $17$ bis $51$, und die
Figure of Merit steigt um Faktoren von $61$ bis $449$, obwohl die Himmelsfläche etwa
$28$-mal kleiner ist. Mit der aktuellen KiDS-Fläche lassen sich die beiden Szenarien jedoch noch nicht
unterscheiden. Eine Vorhersage für den Legacy Survey of Space and Time (LSST) zeigt, dass
kommende Galaxiendurchmusterungen die Aussichten auf eine Unterscheidung zwischen prompten und
verzögerten FRB-Ursprüngen erheblich verbessern werden.
\end{foreignlanguage}